% Manual · Laboratorio 01 — Windows: radiodifusión FM con HackRF One, SDRangel y SIGLENT.

\PortadillaLaboratorio{Laboratorio 01 · Windows 10/11 · 210 minutos}{De las ondas de radio al sonido: recepción FM, SDR y análisis de espectro}{Su misión: sintonizar una emisora de FM real, recuperar su audio y medir la misma señal con un analizador de espectro. Al final sabrá explicar qué le dijo cada instrumento y qué no podía decirle.}

\section{Ficha del laboratorio}
\label{sec:w-plan}

\begin{longtable}{@{}>{\raggedright\arraybackslash}p{0.20\linewidth}>{\raggedright\arraybackslash}p{0.73\linewidth}@{}}
\toprule
\endhead
\textbf{Al terminar podrá} & montar y verificar una estación SDR receptora; configurar SDRangel distinguiendo el dispositivo del canal; localizar, demodular y escuchar una emisora; medirla con el SIGLENT (centro, span, referencia, RBW, VBW, detector, marcadores); comprobar con experimentos el efecto de la sintonía, el ancho de banda y la ganancia; y separar la detección de energía de la recuperación de información. \\
\textbf{Antes de venir} & Responda el cuestionario diagnóstico (Anexo~B, p.~\pageref{B-sec:b-diagw}) y lea el \anexoA{sec:a-fm}. \\
\textbf{Equipo} & HackRF One, antena ANT500, cable de extensión SMA, adaptador N -- SMA, SIGLENT SSA3032X Plus, PC con Windows y auriculares (sección~\ref{sec:banco}). \\
\textbf{Registro} & Hojas W1 a W9 del Anexo B. \\
\bottomrule
\end{longtable}

\begin{longtable}{@{}>{\raggedright\arraybackslash}p{0.62\linewidth}>{\raggedleft\arraybackslash}p{0.12\linewidth}>{\raggedleft\arraybackslash}p{0.16\linewidth}@{}}
\toprule
\textbf{Etapa} & \textbf{Minutos} & \textbf{Acumulado} \\
\midrule
\endhead
Orientación: misión, equipo y reglas & 15 & 15 \\
1 · Montar la estación & 15 & 30 \\
2 · Windows reconoce el HackRF & 15 & 45 \\
3 · El receptor en SDRangel & 25 & 70 \\
4 · Encontrar y escuchar una emisora & 40 & 110 \\
5 · Medir con el analizador & 40 & 150 \\
6 · Cuatro experimentos & 40 & 190 \\
7 · Cierre y entrega & 20 & 210 \\
\bottomrule
\end{longtable}

\begin{figure}[H]
\centering
\resizebox{\linewidth}{!}{% Laboratorio 01: cadena de recepción FM y configuración secuencial del analizador.
\begin{tikzpicture}[node distance=5mm and 7mm]
  % Configuración 1: SDR
  \node[bloqueext,minimum width=23mm] (tx) {Emisora FM\\comercial\\(88--108 MHz)};
  \node[bloque,right=14mm of tx,minimum width=23mm] (ant) {Antena \textbf{ANT500}\\extendida a\\$\approx$ 75 cm};
  \node[bloque,right=of ant,minimum width=19mm] (coax) {Cable de\\extensión SMA\\(si hace falta)};
  \node[bloque,right=of coax,minimum width=26mm] (hrf) {\textbf{HackRF One}\\puerto \textbf{ANTENNA}\\(SMA hembra)};
  \draw[flecha,decorate,decoration={snake,amplitude=1.1mm,segment length=3mm,post length=2.5mm}] (tx) -- (ant);
  \node[rotulo,above=0.5mm] at ($(tx.east)!0.5!(ant.west)+(0,3mm)$) {onda EM};
  \draw[flecha] (ant) -- (coax);
  \draw[flecha] (coax) -- (hrf);

  \node[bloque,below=12mm of hrf,minimum width=26mm] (pc) {PC Windows\\controlador WinUSB};
  \draw[flecha] (hrf) -- node[right,rotulo,align=left]{USB 2.0\\I/Q} (pc);
  \node[bloquealt,left=of pc,minimum width=24mm] (dev) {SDRangel\\\textbf{HackRF Input}\\(frecuencia central)};
  \node[bloquealt,left=of dev,minimum width=26mm] (bfm) {\textbf{Broadcast FM}\\\textbf{Demodulator}\\(desplazamiento $\Delta f$)};
  \node[bloque,left=of bfm,minimum width=17mm] (aud) {Salida\\de audio};
  \draw[flecha] (pc) -- (dev);
  \draw[flecha] (dev) -- (bfm);
  \draw[flecha] (bfm) -- (aud);

  % Configuración 2: analizador, secuencial
  \node[bloque,below=17mm of aud.south west,anchor=north west,minimum width=23mm] (ant2) {La \emph{misma}\\antena ANT500};
  \node[bloque,right=9mm of ant2,minimum width=29mm] (adp) {Adaptador\\N macho --\\SMA hembra};
  \node[bloqueinst,right=9mm of adp,minimum width=36mm] (sa) {\textbf{SIGLENT SSA3032X Plus}\\\textbf{RF INPUT} 50~$\Omega$ (N hembra)\\generador TG \emph{apagado}};
  \draw[flecha] (ant2) -- (adp);
  \draw[flecha] (adp) -- (sa);

  \begin{scope}[on background layer]
    \node[draw=MARule,rounded corners=3pt,fit=(tx)(hrf)(pc)(aud),inner sep=3mm,label={[rotulo,anchor=south west]north west:\textbf{Configuración 1} · recepción y demodulación con el SDR}] (c1) {};
    \node[draw=MARule,rounded corners=3pt,fit=(ant2)(adp)(sa),inner sep=3mm,label={[rotulo,anchor=south west]north west:\textbf{Configuración 2} · después de la 1: la antena se traslada al analizador (no hay divisor)}] {};
  \end{scope}
\end{tikzpicture}
}
\caption{Las dos configuraciones del laboratorio. Primero, la antena en el HackRF (etapas 1--4); después, la misma antena en el analizador (etapa 5). Los dos instrumentos nunca se conectan entre sí.}
\label{fig:cadenafm}
\end{figure}

% ---------------------------------------------------------------------------
\section{Etapa 1 · Montar la estación}
\label{sec:w-equipo}

\Puente{Todo empieza en el conector. Un receptor bien configurado con la antena en el puerto equivocado no recibe nada y no da ningún error.}

\Paso{Identifique los puertos del HackRF}
{Carcasa del HackRF One.}
{Tres conectores SMA hembra: \textbf{ANTENNA} (señal de radio) y \textbf{CLKIN}/\textbf{CLKOUT} (reloj). Un USB y dos pulsadores.}
{La antena solo funciona en ANTENNA.}
{Usted señala el puerto ANTENNA sin dudar.}
{}
{Si la serigrafía no se lee, pregunte antes de conectar.}

\Paso{Monte la antena ANT500}
{Puerto ANTENNA; la tuerca SMA de la antena.}
{Enrósquela a mano. Extiéndala hasta unos \valor{75 cm} y póngala vertical.}
{75 cm es un cuarto de la longitud de onda a 100 MHz: la antena resuena en el centro de la banda de FM.}
{La antena queda firme y vertical.}
{Mueva la antena con suavidad: el conector no tiene holgura.}
{Si la rosca no entra con suavidad, no fuerce: compruebe que no mezcla conectores distintos (\anexoA{sec:a-antenas}).}

\Paso{Conecte el USB}
{Un puerto USB del PC, directo (sin concentrador).}
{Cable de datos del HackRF.}
{El flujo de muestras necesita un enlace USB 2.0 estable.}
{Se encienden los LED del HackRF; Windows emite el sonido de dispositivo conectado.}
{}
{Ningún LED: otro cable y otro puerto.}

Anote los datos del puesto en la \hoja{W1}.

\begin{logro}
La estación está montada: antena en el puerto correcto, a la longitud correcta, y el HackRF alimentado.
\end{logro}

% ---------------------------------------------------------------------------
\section{Etapa 2 · Windows reconoce el HackRF}
\label{sec:w-windows}

\Puente{Antes de abrir SDRangel, compruebe que Windows ve el HackRF. Si este paso falla, ningún ajuste posterior lo arreglará.}

\Paso{Compruebe el controlador}
{\tecla{Win}+\tecla{X} $\rightarrow$ \menu{Administrador de dispositivos}.}
{Busque \textbf{HackRF One} entre los dispositivos USB.}
{Windows necesita el controlador WinUSB para que SDRangel use el HackRF.}
{«HackRF One» sin triángulo amarillo.}
{Desconecte el USB: la entrada desaparece. Reconéctelo: vuelve.}
{Triángulo amarillo en \menu{Otros dispositivos}: avise al instructor (\anexoA{sec:a-windows}).}

\Comando{hackrf_info}
{Pedir al HackRF que se identifique. Ábralo en \menu{Inicio $\rightarrow$ radioconda Prompt}, con SDRangel cerrado.}
{\texttt{Found HackRF}, \texttt{Board ID Number: 2 (HackRF One)}, número de serie y versión de firmware.}
{\texttt{No HackRF boards found}.}
{Cierre cualquier programa SDR, reconecte el USB y repita.}

\begin{logro}
Windows reconoce el HackRF y usted conoce su número de serie y su firmware (\hoja{W1}).
\end{logro}

% ---------------------------------------------------------------------------
\section{Etapa 3 · El receptor en SDRangel}
\label{sec:w-sdrangel}

\Puente{SDRangel separa dos ideas: el \emph{dispositivo} (el HackRF, que observa una ventana de frecuencias) y los \emph{canales} (procesadores que extraen una señal de esa ventana). Primero se configura el dispositivo.}

\Paso{Abra SDRangel}
{\menu{Inicio $\rightarrow$ SDRangel}. Cierre antes el radioconda Prompt.}
{Si no ve la barra de herramientas: \menu{Workspaces $\rightarrow$ New}.}
{El HackRF admite una sola aplicación a la vez.}
{La ventana principal con un espacio de trabajo vacío.}
{}
{}

\Paso{Añada el receptor HackRF}
{\tecla{Ctrl}+\tecla{R} $\rightarrow$ diálogo \emph{Choose a receiving device\ldots}}
{La entrada \textbf{HackRF} con su número de serie $\rightarrow$ \gui{OK}.}
{}
{Una ventana \textbf{HackRF Input} y una ventana de espectro.}
{El número de serie coincide con el de la etapa 2.}
{Si solo aparecen \texttt{FileInput}, \texttt{TestSource} o similares, SDRangel no ve el HackRF: ciérrelo, reconecte el USB y repita.}
\label{paso:w-seldisp}

\Paso{Configure el dispositivo}
{Ventana \textbf{HackRF Input}.}
{\gui{Bias tee} y \gui{RF amp} \valor{desactivados}; selector \gui{SR}; tasa \valor{2\,400\,000 S/s}; decimación \valor{1}; \gui{Fc} \valor{Cen}; \gui{DC} e \gui{IQ} \valor{activados}; \gui{BBF} \valor{automático}; \gui{LNA} \valor{16 dB}; \gui{VGA} \valor{16 dB}.}
{Ganancias moderadas para no saturar con emisoras fuertes; 2,4 MSa/s muestran 2,4 MHz de espectro a la vez.}
{Los valores quedan como se indican.}
{Léalos en voz alta con su compañero.}
{}

\Paso{Arranque en el inicio de la banda}
{Dial de frecuencia (en kHz) y botón de arranque.}
{\valor{89\,000 kHz}. El triángulo azul pasa a cuadrado verde.}
{}
{Un piso de ruido con picos; la cascada se desplaza hacia abajo.}
{Esto es solo energía (nivel 1): todavía no ha identificado ninguna emisora.}
{Cuadrado rojo: el HackRF se desconectó o lo usa otro programa. Detenga, reconecte y arranque.}

\Profundice{cómo se convierte la onda en muestras I/Q y qué hace cada ganancia, \anexoA{sec:a-sdr}; espectro y cascada, \anexoA{sec:a-espectro}}

\begin{logro}
El HackRF entrega muestras y usted ve 2,4 MHz de la banda de FM en tiempo real.
\end{logro}

% ---------------------------------------------------------------------------
\section{Etapa 4 · Encontrar y escuchar una emisora}
\label{sec:w-fm}

\Puente{Ahora se pasa de la energía a la información: elegir una emisora, situar un canal sobre ella y escuchar. Escriba primero su predicción en la \hoja{W2}.}

\Paso{Recorra la banda}
{Dial de frecuencia del dispositivo.}
{89, 91, 93, \ldots, 107 MHz; espere unos segundos en cada posición.}
{Con una ventana de 2,4 MHz, saltos de 2 MHz cubren la banda con solape.}
{Bloques de unos 200 kHz que sobresalen del ruido y forman trazas continuas en la cascada.}
{Anote las 3--5 emisoras más altas en la \hoja{W2}. Ignore la línea estrecha que aparece \emph{siempre} en el centro: es un artefacto del receptor (pico DC), no una emisora.}
{Sin emisoras: revise el conector y lleve la antena a una ventana con el cable de extensión.}

\begin{figure}[H]
\centering
\resizebox{0.9\linewidth}{!}{% Frecuencia del dispositivo frente a desplazamiento del canal (sintonía con offset).
\begin{tikzpicture}[x=1mm,y=1mm,font=\sffamily\scriptsize]
  \fill[MASurfaceAlt] (10,0) rectangle (130,30);
  \draw[MAMuted] (10,0) rectangle (130,30);
  \node[rotulo,anchor=south] at (70,31) {ventana observada = tasa de muestreo compleja (2,4 MSa/s $\Rightarrow$ 2,4 MHz)};
  \draw[flecha] (4,0) -- (138,0) node[right,rotulo]{$f$};
  \draw[MAMuted] (10,4) -- (130,4);
  % pico DC en el centro del dispositivo
  \draw[MADanger,line width=1.4pt] (70,4) -- (70,20);
  \node[MADanger,align=center,anchor=south] at (70,20.5) {pico DC\\(artefacto)};
  % emisora a -300 kHz
  \draw[line width=0.9pt,color=MAPrimaryDark,fill=MAPrimarySoft]
     plot[smooth] coordinates {(48,4.2) (50,8) (52,13) (55,16) (58,13) (60,8) (62,4.2)};
  % canal
  \draw[MASecondary,line width=0.8pt] (48,2) rectangle (62,25);
  \node[MASecondary,anchor=east,align=right] at (47,20) {canal BFM\\(180 kHz)};
  % marcas de frecuencia
  \draw[dashed] (55,0) -- (55,-2) node[below]{$f_{\mathrm{emisora}}$};
  \draw[dashed] (70,0) -- (70,-2) node[below]{$f_{\mathrm{dispositivo}}$};
  \draw[<->,MASecondary,line width=0.8pt] (55,-9) -- (70,-9);
  \node[MASecondary,anchor=north] at (62.5,-10) {$\Delta f = f_{\mathrm{emisora}} - f_{\mathrm{dispositivo}}$ (aquí $-300$ kHz)};
  \node[rotulo,anchor=west,align=left] at (86,15) {Sintonice el dispositivo \textbf{300 kHz por encima}\\de la emisora y lleve el canal a $-300$ kHz:\\así no demodula sobre el pico DC.};
\end{tikzpicture}
}
\caption{El dispositivo fija el centro de la ventana; el canal elige qué parte de ella se demodula. Centrar el dispositivo 300 kHz por encima de la emisora la aleja del pico DC.}
\label{fig:offset}
\end{figure}

\Paso{Sitúe la emisora fuera del centro}
{Dial de frecuencia del dispositivo.}
{$f_{\mathrm{dispositivo}} = f_{\mathrm{emisora}} + 300$ kHz. Ejemplo: emisora en 100,9 MHz $\Rightarrow$ \valor{101\,200 kHz}.}
{Demodular encima del pico DC degrada el audio.}
{La emisora aparece 300 kHz a la izquierda del centro.}
{}
{}

\Paso{Añada el demodulador de FM de radiodifusión}
{Icono de añadir canal, arriba en la ventana \textbf{HackRF Input}.}
{\textbf{Broadcast FM Demodulator} $\rightarrow$ \gui{Apply} $\rightarrow$ \gui{Close}.}
{Es el demodulador para FM comercial (canal ancho, estéreo y RDS); \emph{NFM} es para radios de voz.}
{Una ventana \textbf{Broadcast FM Demod} y una franja de color sobre el espectro.}
{}
{}

\Paso{Configure el canal}
{Ventana \textbf{Broadcast FM Demod}.}
{Desplazamiento \valor{$-300$ kHz} ($-300\,000$ Hz en los diales: rueda del ratón o flechas sobre cada dígito); \gui{RF BW} \valor{180 kHz}; \gui{AF BW} \valor{15 kHz}; de-énfasis \valor{75 $\mu$s}; volumen \valor{2--5}; squelch al mínimo.}
{El desplazamiento lleva el canal a la emisora. En América el de-énfasis es de 75 $\mu$s; SDRangel trae 50 $\mu$s por defecto.}
{La franja cubre la emisora y el medidor del canal se mueve.}
{La potencia del canal sube varios dB sobre la emisora y baja al apartarlo.}
{La franja no cubre la emisora: revise el signo del desplazamiento.}

\Paso{Escuche}
{Auriculares; ventana \textbf{Broadcast FM Demod}.}
{Audio sin silenciar.}
{El audio inteligible demuestra la demodulación (nivel 2).}
{Voz o música reconocibles. Si activa \emph{Stereo} y la señal es buena, el botón se pone verde; con \emph{RDS} puede aparecer el nombre de la emisora.}
{Aparte el canal 200 kHz: el audio se pierde o cambia de emisora. Vuelva a $-300$ kHz.}
{Hay señal pero no se oye nada: mute, squelch y \menu{Preferences $\rightarrow$ Audio} (\anexoA{sec:a-diagfm}).}

Complete la \hoja{W3}.

\Profundice{por qué una emisora de FM ocupa unos 200 kHz y qué son la señal MPX, el piloto y el RDS, \anexoA{sec:a-fm}}

\begin{logro}
Ha recuperado información de una señal de radio real: nivel 2 de la escalera de evidencia.
\end{logro}

% ---------------------------------------------------------------------------
\section{Etapa 5 · Medir con el analizador}
\label{sec:w-siglent}

\Puente{El SDR le permitió escuchar. El analizador le permitirá \emph{medir}: niveles absolutos en dBm y anchos con marcadores. Se usa la misma antena, trasladada, para comparar en igualdad de condiciones.}

\Paso{Prepare el analizador}
{Interruptor de encendido; tecla \gui{Preset}; menú \gui{TG}.}
{\gui{Preset} tras el arranque. Generador de seguimiento (TG) \valor{apagado}.}
{\gui{Preset} parte de un estado conocido; el TG es una salida y debe quedar apagado.}
{Barrido completo con el piso de ruido.}
{Ningún indicador de TG activo en pantalla.}
{}

\Paso{Traslade la antena}
{Detenga el HackRF en SDRangel. ANT500 $\rightarrow$ adaptador N macho -- SMA hembra $\rightarrow$ \textbf{RF INPUT}.}
{Toque una superficie metálica antes de conectar.}
{No hay divisor: la antena se usa por turnos.}
{La antena queda firme y vertical sobre el analizador.}
{}
{}

\Paso{Vea la banda completa}
{\gui{Frequency $\rightarrow$ Center Freq}; \gui{Span}; \gui{Amplitude $\rightarrow$ Ref Level}; \gui{Amplitude $\rightarrow$ Attenuator}; \gui{BW}.}
{Centro \valor{98 MHz}; span \valor{20 MHz}; referencia \valor{$-20$ dBm}; atenuación \valor{Auto}; preamplificador \valor{apagado}; RBW y VBW \valor{Auto}.}
{20 MHz cubren 88--108 MHz.}
{Varias emisoras como picos sobre el ruido.}
{Las más altas coinciden con su \hoja{W2}.}
{Solo ruido: active \gui{Amplitude $\rightarrow$ Preamp} y baje la referencia a $-40$ dBm.}

\Paso{Acérquese a su emisora}
{\gui{Frequency}, \gui{Span}, \gui{BW}, \gui{Detect}.}
{Centro $= f_{\mathrm{emisora}}$; span \valor{1 MHz}; RBW \valor{10 kHz}; VBW \valor{10 kHz}; detector \valor{Pos Peak}.}
{Una RBW de 10 kHz resuelve la forma de un canal de 200 kHz.}
{Un bloque de unos 200 kHz cuya forma cambia con el audio.}
{}
{}

\Paso{Acumule y mida}
{\gui{Trace} $\rightarrow$ \valor{Max Hold}; después \gui{Peak} y \gui{Marker Fn $\rightarrow$ N dB BW}.}
{Max Hold durante 30 s; marcador en el pico; \valor{N = $-26$ dB}.}
{El máximo acumulado muestra el ancho realmente ocupado; \emph{N dB BW} mide el ancho entre los puntos a $N$ dB del pico.}
{Frecuencia y nivel del pico en dBm y un ancho del orden de 150--300 kHz.}
{}
{Sin \emph{N dB BW}: dos marcadores delta a $-26$ dB a cada lado del pico.}

Complete la \hoja{W4}. Si le queda tiempo, pruebe la tecla \gui{Demod} del analizador: también puede demodular FM y reproducirla por su conector de auriculares.

\begin{importante}[Qué significa ese número en dBm]
Es la potencia \emph{en la entrada del analizador}, medida con su RBW, a través de esta antena y este adaptador. No es la potencia del transmisor ni puede compararse directamente con los dB de SDRangel, que son relativos a la escala digital (\anexoA{sec:a-analizador}).
\end{importante}

Al terminar, devuelva la antena al HackRF y arranque de nuevo el receptor.

\begin{logro}
Tiene una medida absoluta (dBm) y un ancho ocupado de la misma emisora que escuchó.
\end{logro}

% ---------------------------------------------------------------------------
\section{Etapa 6 · Cuatro experimentos}
\label{sec:w-exp}

\Puente{Ya sabe obtener el resultado; ahora va a averiguar de qué depende. En cada experimento escriba primero su predicción en la hoja indicada, cambie \emph{una sola} cosa, observe y explique.}

\begin{longtable}{@{}>{\raggedright\arraybackslash}p{0.17\linewidth}>{\raggedright\arraybackslash}p{0.55\linewidth}>{\raggedright\arraybackslash}p{0.20\linewidth}@{}}
\toprule
\textbf{Experimento} & \textbf{Qué cambia} & \textbf{Registro} \\
\midrule
\endhead
1 · Sintonía & Desplazamiento del canal: $-250$, $-200$ y $-100$ kHz (50, 100 y 200 kHz fuera de la emisora). Vuelva a $-300$. & \hoja{W5} \\
2 · Ancho de banda & \gui{RF BW} 80, 180 y 250 kHz; después \gui{AF BW} 3, 8 y 15 kHz. Termine en 180/15. & \hoja{W6} \\
3 · Ganancia & LNA/VGA: 8/8, 16/16, 32/30 y 40/40 dB, esperando 10 s en cada una. \gui{RF amp} apagado. Termine en 16/16. & \hoja{W7} \\
4 · Instrumentos & Compare lo que puede y no puede decirle cada instrumento. & \hoja{W8} \\
\bottomrule
\end{longtable}

\Profundice{cómo reconocer la saturación y por qué la ganancia no mejora siempre la señal, \anexoA{sec:a-sdr}}

\begin{logro}
Puede explicar, con sus propios datos, qué controla la calidad de la recepción y qué no.
\end{logro}

% ---------------------------------------------------------------------------
\section{Etapa 7 · Cierre y entrega}
\label{sec:w-evid}

\begin{enumerate}
  \item Detenga el HackRF en SDRangel y cierre el programa; desconecte el USB y después la antena.
  \item Apague el analizador y guarde el adaptador N.
  \item Complete la interpretación final (\hoja{W9}) y reúna los entregables de la lista del Anexo B (p.~\pageref{B-sec:b-entregaw}).
  \item Responda el cuestionario final (Anexo B, p.~\pageref{B-sec:b-finalw}).
\end{enumerate}

\begin{logro}[Laboratorio 01 completado]
Ha recorrido toda la cadena de un receptor, de la onda al sonido, y ha medido la señal con un instrumento independiente. En el Laboratorio 02 aplicará lo mismo a señales mucho más difíciles: ráfagas de 120 $\mu$s que solo tienen sentido si se decodifican.
\end{logro}
